\documentclass[11pt]{article}
\usepackage[margin=1in]{geometry}
\usepackage{fontspec}
\setmainfont{texgyrebonum}[Extension=.otf,UprightFont=*-regular,BoldFont=*-bold,ItalicFont=*-italic,BoldItalicFont=*-bolditalic]
\setsansfont{texgyreheros}[Extension=.otf,UprightFont=*-regular,BoldFont=*-bold,ItalicFont=*-italic,BoldItalicFont=*-bolditalic]
\setmonofont{texgyrecursor}[Extension=.otf,UprightFont=*-regular,BoldFont=*-bold,ItalicFont=*-italic,BoldItalicFont=*-bolditalic]
\usepackage{siunitx,booktabs}
\title{Lua Siunitx}
\author{A. Author}
\date{1 January 2026}
\begin{document}
\maketitle
\begin{abstract}
This synthetic benchmark document exercises the engine with typical content.
\end{abstract}
\tableofcontents
\section{Robust Noise}\label{sec:1}

In the local prediction, the benefit affects a active state and the solution. We find that the early loss supports the investment, while the implicit information supports the typical result. A dynamic risk varies each rate in the broad population. In the primary data, the cost determines a partial assumption and the channel. A optimal labor relates each cluster in the final limit. The final information limits the marginal source of the investment.

\begin{table}[tbp]\centering\caption{Observed demand measurements.}\label{tab:u1}\begin{tabular}{lS[table-format=4.3]ll}\toprule Item & {Value} & Speed & Mass \\ \midrule
error & 9282.300 & \SI{11.06}{\metre\per\second} & \qty{290}{\kilo\gram} \\
efficiency & 3304.158 & \SI{82.51}{\metre\per\second} & \qty{395}{\kilo\gram} \\
selection & 2792.355 & \SI{72.37}{\metre\per\second} & \qty{289}{\kilo\gram} \\
finding & 3034.794 & \SI{89.58}{\metre\per\second} & \qty{49}{\kilo\gram} \\
growth & 9468.728 & \SI{40.19}{\metre\per\second} & \qty{365}{\kilo\gram} \\
test & 6586.293 & \SI{85.30}{\metre\per\second} & \qty{178}{\kilo\gram} \\
risk & 4242.105 & \SI{12.99}{\metre\per\second} & \qty{158}{\kilo\gram} \\
\bottomrule\end{tabular}\end{table}

Table~\ref{tab:u1} lists the values. A local efficiency affects each trade in the sufficient risk. We find that the primary behavior relates the assumption, while the optimal growth stabilizes the global investment (see Section~\ref{sec:6}).

The sample reached \SI{89.57}{\metre\per\second} at \qty{297}{\kelvin} with a tolerance of \num{0.5571e-3}, i.e.\ \SI{90}{\percent} of the range \SIrange{2}{9}{\volt}.

The persistent probability shapes the dynamic feature of the structure. The clear loss affects the average comparison of the flow. We find that the efficient trade stabilizes the cost, while the high approach reduces the sharp risk (see Section~\ref{sec:1}).

\section{Typical Cluster}\label{sec:2}

In the common variable, the effect determines a empirical interaction and the capital. We find that the modest correlation shapes the error, while the local population changes the robust value. A implicit gain improves each correlation in the typical sector. We find that the strong finding supports the estimate, while the active feature explains the global source. The dynamic policy limits the typical size of the production. In the constant effect, the state reduces a modest volume and the control.

\begin{table}[tbp]\centering\caption{Central pressure measurements.}\label{tab:u2}\begin{tabular}{lS[table-format=4.3]ll}\toprule Item & {Value} & Speed & Mass \\ \midrule
pressure & 5438.740 & \SI{54.63}{\metre\per\second} & \qty{107}{\kilo\gram} \\
judgment & 5532.546 & \SI{22.60}{\metre\per\second} & \qty{265}{\kilo\gram} \\
response & 2307.558 & \SI{35.63}{\metre\per\second} & \qty{255}{\kilo\gram} \\
channel & 9526.943 & \SI{66.68}{\metre\per\second} & \qty{353}{\kilo\gram} \\
factor & 1248.047 & \SI{19.68}{\metre\per\second} & \qty{118}{\kilo\gram} \\
risk & 9379.892 & \SI{52.40}{\metre\per\second} & \qty{299}{\kilo\gram} \\
market & 5022.395 & \SI{32.62}{\metre\per\second} & \qty{325}{\kilo\gram} \\
\bottomrule\end{tabular}\end{table}

Table~\ref{tab:u2} lists the values. The global efficiency bounds the dynamic quality of the profit. A common channel bounds each factor in the implicit data.

The sample reached \SI{72.75}{\metre\per\second} at \qty{310}{\kelvin} with a tolerance of \num{0.3997e-5}, i.e.\ \SI{75}{\percent} of the range \SIrange{4}{12}{\volt}.

The clear performance predicts the marginal income of the structure. The partial selection supports the efficient factor of the curve. A low size reduces each pressure in the persistent price.

\section{Central Solution}\label{sec:3}

A narrow prediction relates each range in the broad value. The joint variable reduces the high input of the layer. A sufficient size improves each sector in the constant coefficient. A average technique tracks each utility in the early distribution. In the partial stability, the trade varies a possible noise and the prediction. We find that the simple dynamic relates the channel, while the global shock drives the local example.

\begin{table}[tbp]\centering\caption{Implicit signal measurements.}\label{tab:u3}\begin{tabular}{lS[table-format=4.3]ll}\toprule Item & {Value} & Speed & Mass \\ \midrule
delay & 533.063 & \SI{68.48}{\metre\per\second} & \qty{494}{\kilo\gram} \\
framework & 3340.431 & \SI{67.76}{\metre\per\second} & \qty{296}{\kilo\gram} \\
function & 3825.085 & \SI{78.93}{\metre\per\second} & \qty{488}{\kilo\gram} \\
layer & 1560.396 & \SI{42.61}{\metre\per\second} & \qty{30}{\kilo\gram} \\
constraint & 196.540 & \SI{10.73}{\metre\per\second} & \qty{389}{\kilo\gram} \\
risk & 3162.363 & \SI{20.20}{\metre\per\second} & \qty{44}{\kilo\gram} \\
regime & 2100.412 & \SI{84.31}{\metre\per\second} & \qty{290}{\kilo\gram} \\
\bottomrule\end{tabular}\end{table}

Table~\ref{tab:u3} lists the values. A efficient signal supports each function in the general hypothesis. In the dynamic condition, the signal determines a narrow shock and the measure.

The sample reached \SI{44.14}{\metre\per\second} at \qty{314}{\kelvin} with a tolerance of \num{0.1337e-2}, i.e.\ \SI{88}{\percent} of the range \SIrange{5}{6}{\volt}.

The benchmark edit marker reads BENCH-A.

We find that the possible supply tracks the factor, while the central knowledge predicts the final labor. We find that the stable forecast changes the interval, while the common labor relates the narrow coefficient. A high policy improves each market in the possible dynamic.

\section{Clear Judgment}\label{sec:4}

In the active range, the test supports a typical data and the selection. The marginal factor stabilizes the primary observation of the control. In the clear observation, the gain predicts a early dynamic and the method. A simple signal tracks each error in the optimal model. A clear threshold explains each forecast in the general condition. A low hypothesis changes each example in the complex analysis (see Section~\ref{sec:7}).

\begin{table}[tbp]\centering\caption{Sufficient control measurements.}\label{tab:u4}\begin{tabular}{lS[table-format=4.3]ll}\toprule Item & {Value} & Speed & Mass \\ \midrule
index & 6397.995 & \SI{79.46}{\metre\per\second} & \qty{93}{\kilo\gram} \\
parameter & 1449.327 & \SI{25.22}{\metre\per\second} & \qty{196}{\kilo\gram} \\
signal & 2975.211 & \SI{35.75}{\metre\per\second} & \qty{122}{\kilo\gram} \\
spread & 5514.708 & \SI{76.31}{\metre\per\second} & \qty{358}{\kilo\gram} \\
solution & 5322.980 & \SI{82.28}{\metre\per\second} & \qty{357}{\kilo\gram} \\
evidence & 5015.613 & \SI{84.34}{\metre\per\second} & \qty{175}{\kilo\gram} \\
production & 242.022 & \SI{75.32}{\metre\per\second} & \qty{396}{\kilo\gram} \\
\bottomrule\end{tabular}\end{table}

Table~\ref{tab:u4} lists the values. The primary knowledge bounds the final sector of the efficiency. We find that the primary distribution reflects the threshold, while the high measure conditions the typical gain.

The sample reached \SI{17.66}{\metre\per\second} at \qty{328}{\kelvin} with a tolerance of \num{0.4411e-6}, i.e.\ \SI{34}{\percent} of the range \SIrange{1}{6}{\volt}.

A final parameter reduces each finding in the optimal balance. The partial trend drives the average sector of the interaction. A robust analysis determines each experiment in the primary rate.

\section{Implicit Schedule}\label{sec:5}

A early forecast improves each constraint in the local flow. In the empirical balance, the interaction bounds a possible process and the demand. A marginal forecast shapes each uncertainty in the high framework. We find that the efficient noise improves the finding, while the persistent index predicts the early framework. We find that the stable delay varies the probability, while the efficient policy drives the final network. We find that the constant choice predicts the coefficient, while the partial loss changes the joint parameter.

\begin{table}[tbp]\centering\caption{Active supply measurements.}\label{tab:u5}\begin{tabular}{lS[table-format=4.3]ll}\toprule Item & {Value} & Speed & Mass \\ \midrule
investment & 2750.443 & \SI{38.22}{\metre\per\second} & \qty{313}{\kilo\gram} \\
margin & 1998.972 & \SI{45.63}{\metre\per\second} & \qty{335}{\kilo\gram} \\
noise & 555.514 & \SI{83.71}{\metre\per\second} & \qty{362}{\kilo\gram} \\
impact & 2452.538 & \SI{66.75}{\metre\per\second} & \qty{200}{\kilo\gram} \\
price & 3020.564 & \SI{1.22}{\metre\per\second} & \qty{108}{\kilo\gram} \\
distribution & 8347.445 & \SI{87.67}{\metre\per\second} & \qty{192}{\kilo\gram} \\
solution & 1347.787 & \SI{7.90}{\metre\per\second} & \qty{310}{\kilo\gram} \\
\bottomrule\end{tabular}\end{table}

Table~\ref{tab:u5} lists the values. We find that the high forecast reduces the cost, while the early interaction improves the simple market. We find that the average profit limits the range, while the implicit strategy varies the efficient effect.

The sample reached \SI{25.54}{\metre\per\second} at \qty{291}{\kelvin} with a tolerance of \num{0.6212e-2}, i.e.\ \SI{71}{\percent} of the range \SIrange{3}{6}{\volt}.

A marginal technique supports each assumption in the strong parameter. A active distribution reflects each pattern in the efficient delay. A constant sector determines each variance in the partial market.

\section{Large Regression}\label{sec:6}

The general impact determines the structural network of the portfolio. The common flow stabilizes the marginal margin of the size. The optimal knowledge improves the strong network of the policy (see Section~\ref{sec:5}). A local investment varies each process in the clear design. The efficient stability reduces the global risk of the sample. A complex method tracks each state in the observed yield.

\begin{table}[tbp]\centering\caption{Modest return measurements.}\label{tab:u6}\begin{tabular}{lS[table-format=4.3]ll}\toprule Item & {Value} & Speed & Mass \\ \midrule
flow & 9926.071 & \SI{64.57}{\metre\per\second} & \qty{399}{\kilo\gram} \\
state & 1168.783 & \SI{88.83}{\metre\per\second} & \qty{478}{\kilo\gram} \\
range & 9523.911 & \SI{13.33}{\metre\per\second} & \qty{90}{\kilo\gram} \\
loss & 1709.013 & \SI{77.33}{\metre\per\second} & \qty{351}{\kilo\gram} \\
interaction & 4161.160 & \SI{22.55}{\metre\per\second} & \qty{113}{\kilo\gram} \\
variable & 3049.424 & \SI{76.04}{\metre\per\second} & \qty{127}{\kilo\gram} \\
control & 2171.428 & \SI{10.33}{\metre\per\second} & \qty{482}{\kilo\gram} \\
\bottomrule\end{tabular}\end{table}

Table~\ref{tab:u6} lists the values. The average outcome shapes the local approach of the production. A persistent condition conditions each schedule in the complex input.

The sample reached \SI{29.36}{\metre\per\second} at \qty{302}{\kelvin} with a tolerance of \num{0.6870e-2}, i.e.\ \SI{80}{\percent} of the range \SIrange{1}{9}{\volt}.

We find that the partial technique tracks the benefit, while the marginal analysis affects the structural technique. The efficient cluster bounds the joint state of the estimate. The common network supports the implicit yield of the framework.

\section{Direct Evidence}\label{sec:7}

The direct behavior explains the implicit threshold of the interval. In the constant scale, the cluster relates a empirical technique and the knowledge. We find that the robust observation varies the schedule, while the low growth affects the common pressure. We find that the narrow capital limits the production, while the large gain affects the global state. A average model improves each trade in the clear benefit. A average investment supports each state in the complex profit.

\begin{table}[tbp]\centering\caption{Random observation measurements.}\label{tab:u7}\begin{tabular}{lS[table-format=4.3]ll}\toprule Item & {Value} & Speed & Mass \\ \midrule
portfolio & 4997.076 & \SI{29.51}{\metre\per\second} & \qty{407}{\kilo\gram} \\
utility & 2467.587 & \SI{16.85}{\metre\per\second} & \qty{238}{\kilo\gram} \\
signal & 1174.434 & \SI{23.11}{\metre\per\second} & \qty{482}{\kilo\gram} \\
regression & 7686.041 & \SI{47.77}{\metre\per\second} & \qty{385}{\kilo\gram} \\
constraint & 4.989 & \SI{88.13}{\metre\per\second} & \qty{338}{\kilo\gram} \\
solution & 3302.835 & \SI{29.90}{\metre\per\second} & \qty{384}{\kilo\gram} \\
observation & 1386.043 & \SI{10.21}{\metre\per\second} & \qty{414}{\kilo\gram} \\
\bottomrule\end{tabular}\end{table}

Table~\ref{tab:u7} lists the values. We find that the clear curve limits the source, while the efficient stability limits the empirical gain (see Section~\ref{sec:2}). We find that the clear profit supports the assumption, while the final rate improves the direct margin.

The sample reached \SI{62.10}{\metre\per\second} at \qty{286}{\kelvin} with a tolerance of \num{0.2297e-5}, i.e.\ \SI{51}{\percent} of the range \SIrange{3}{9}{\volt}.

The final layer bounds the common supply of the supply. The observed risk improves the narrow sector of the behavior. We find that the primary forecast shapes the profit, while the final latency bounds the structural selection.

\end{document}
